\documentclass[12pt,a4paper]{article}
\usepackage[T1]{fontenc}
\usepackage[utf8]{inputenc}
\usepackage[polish]{babel}
\usepackage{geometry}
\usepackage{xcolor}
\usepackage{booktabs}
\usepackage{enumitem}
\usepackage{titlesec}
\usepackage{parskip}
\usepackage{microtype}
\usepackage{hyperref}
\usepackage{fancyhdr}

\geometry{margin=2.5cm}

\definecolor{accent}{RGB}{0,80,160}
\titleformat{\section}{\large\bfseries\color{accent}}{\thesection.\ }{0.4em}{}
\titleformat{\subsection}{\normalsize\bfseries}{\thesubsection.\ }{0.4em}{}

\hypersetup{colorlinks=true, linkcolor=accent, urlcolor=accent, citecolor=accent}

\pagestyle{fancy}
\fancyhf{}
\lhead{Obliczenia Superkomputerowe}
\rhead{Przegląd literatury}
\cfoot{\thepage}

\title{\textbf{Przegląd literatury i~istniejących rozwiązań}\\[0.4em]
\large Klasyfikacja aktywności człowieka na~sygnałach IMU\\
z~wykorzystaniem własnej biblioteki sieci neuronowych przyspieszanej na~GPU}
\author{Stanisław Krupa 254977 \and Julia Budziszewska 267769}
\date{Politechnika Wrocławska \\ przedmiot: \textit{Obliczenia Superkomputerowe}}

\begin{document}
\maketitle

\begin{abstract}
\noindent
Niniejsza notatka stanowi przegląd literatury i~istniejących rozwiązań. Celem jest ustalenie tematu i~zakresu projektu, zestawienie metod stosowanych w~rozpoznawaniu aktywności człowieka (HAR) oraz w~obliczeniach na~akceleratorach GPU, a~także uzasadnienie wyboru architektury i~narzędzi obliczeniowych. Na podstawie źródeł przyjmujemy, że rdzeniem projektu będzie własna, wektoryzowana biblioteka sieci neuronowych (NumPy/CuPy) z~jednowymiarową siecią splotową, trenowaną na trzech reprezentacjach wejścia: surowych osiach (\textit{raw}), modułach wektorów (\textit{magnitude}) oraz ich połączeniu (\textit{hybrid}).
\end{abstract}

\section{Temat i~zakres projektu}

\paragraph{Temat.}
Implementacja od podstaw biblioteki głębokiego uczenia o~wspólnym interfejsie CPU/GPU oraz jej zastosowanie do klasyfikacji 18 aktywności człowieka na zbiorze sygnałów inercjalnych KU-HAR.

\paragraph{Uzasadnienie tematu w~kontekście przedmiotu.}
Głębokie uczenie na szeregach czasowych IMU jest zadaniem o~dużej intensywności obliczeniowej: tensory wejścia mają rząd $N \times C \times T$ (tysiące próbek, kilka kanałów, do $1500$ kroków czasowych), a~trening wymaga wielokrotnego splotu, normalizacji wsadowej i~propagacji wstecz. Jest to więc naturalny poligon dla technik obliczeń superkomputerowych: wektoryzacji, wykorzystania SIMT na GPU, porównania przepustowości CPU i~akceleratora oraz oceny narzutu własnej implementacji względem dojrzałych bibliotek (cuBLAS/cuDNN przez PyTorch).

\paragraph{Zakres planowanych prac.}
\begin{enumerate}[leftmargin=1.4em, itemsep=0.25em]
  \item \textbf{Warstwa obliczeniowa:} backend tablicowy z~przełączaniem NumPy~$\leftrightarrow$~CuPy, bez przepisywania kodu warstw.
  \item \textbf{Biblioteka NN:} warstwy \texttt{Dense}, \texttt{Conv1D}, \texttt{BatchNormalization}, aktywacje, dropout, pooling z~maską długości, optymalizatory SGD/Adam, straty MSE i~entropia krzyżowa, pętla \texttt{fit}/\texttt{predict} oraz serializacja modelu.
  \item \textbf{Zadanie aplikacyjne:} zbiór KU-HAR (akcelerometr + żyroskop smartfona, 18 klas, sekwencje o~zmiennej długości).
  \item \textbf{Trzy warianty cech:} \textit{raw} (6 kanałów), \textit{magnitude} (2 kanały -- normy $L_2$), \textit{hybrid} (8 kanałów).
  \item \textbf{Weryfikacja:} ta sama architektura w~PyTorch; punkty odniesienia GRU oraz PatchTST; pomiar czasu CPU vs.\ GPU.
\end{enumerate}
Poza zakresem pozostawiamy: uczenie rozproszone na wielu węzłach, mieszaną precyzję (\texttt{float16}) oraz strojenie architektury metodą przeszukiwania siatki -- ze względu na ograniczony czas przedmiotu.

\section{Istniejące rozwiązania i~źródła}

Poniżej zestawiono pozycje pogrupowane tematycznie. Przy każdym źródle podano \textbf{wniosek} istotny dla projektu.

\subsection{Zbiory danych i~klasyczne HAR}

\begin{description}[style=nextline, leftmargin=0.6em, font=\normalfont]
  \item[\cite{sikder2021kuhar} --- Sikder, Nahid, \emph{KU-HAR}, Pattern Recognit.\ Lett., 2021]
  Autorzy udostępniają heterogeniczny zbiór 18 aktywności z~90 uczestników, rejestrowany wbudowanymi czujnikami smartfona. Jako bazę podają m.in.\ las losowy z~dokładnością rzędu $90\%$ na krótkich podpróbkach. \textbf{Wniosek:} zbiór jest trudniejszy niż klasyczny UCI-HAR (6 klas, ustalone okno), zawiera zarówno stany statyczne (siad, leżenie), jak i~dynamiczne (bieg, pompki), oraz zmienną długość nagrania -- wymaga maskowania paddingu. \textbf{Wybór:} KU-HAR jako zadanie docelowe, nie UCI-HAR, bo lepiej pokazuje skalę obliczeń (długie sekwencje, 18 klas).

  \item[\cite{anguita2013har} --- Anguita i~in., UCI HAR, ESANN 2013]
  Kanoniczny benchmark: cechy ręcznie wyliczone z~okna + SVM. \textbf{Wniosek:} tradycyjne HAR opiera się na inżynierskich statystykach (średnia, energia, FFT). Podejście to jest tanie obliczeniowo, ale słabo przenosi się na nowe aktywności. \textbf{Wybór:} nie budujemy klasyfikatora na cechach ręcznych jako głównego modelu; cechy głębokie mają być uczone ze sygnału.

  \item[\cite{wang2019survey} --- Wang i~in., \emph{Deep learning for sensor-based activity recognition}, 2019]
  Przegląd pokazuje dominację CNN 1D (kanały osi = kanały ``obrazu'') oraz RNN/LSTM dla zależności długozasięgowych. \textbf{Wniosek:} Conv1D na osiach czasowych to standard \emph{data-driven} i~jest dobrze wektoryzowalny. \textbf{Wybór:} CNN 1D jako model główny.
\end{description}

\subsection{Sieci splotowe, rekurencyjne i~transformery}

\begin{description}[style=nextline, leftmargin=0.6em, font=\normalfont]
  \item[\cite{yang2015cnn} --- Yang i~in., IJCAI 2015]
  CNN na wielokanałowych szeregach IMU: splot i~pooling wzdłuż czasu, wspólne wagi między czujnikami. Autorzy podkreślają, że CNN jest wystarczająco szybka do HAR online. \textbf{Wniosek:} lokalne wzorce (kroki, uderzenia) są dobrze opisywane przez filtry 1D o~rosnącym polu odbiorczym. \textbf{Wybór:} trzy warstwy Conv1D z~malejącym jądrem ($31\to 15\to 7$) i~rosnącą liczbą filtrów ($32\to 64\to 128$), stride~$2$ w~głębszych warstwach jako tani substytut poolingu.

  \item[\cite{ronao2016cnn} --- Ronao, Cho, 2016]
  CNN na surowym sygnale smartfona bez ręcznych cech; głębokość i~regularyzacja silnie wpływają na uogólnienie. \textbf{Wniosek:} surowe osie wystarczą, jeśli sieć jest regularizowana. \textbf{Wybór:} BatchNorm, Dropout $0{,}2$ po blokach splotu i~$0{,}5$ przed klasyfikatorem.

  \item[\cite{hammerla2016deep} --- Hammerla i~in., IJCAI 2016]
  Porównanie CNN, LSTM i~modeli hybrydowych na noszonych czujnikach. LSTM bywa lepszy przy długim kontekście, ale trening jest sekwencyjny. \textbf{Wniosek:} RNN gorzej wykorzystuje równoległość GPU niż splot. \textbf{Wybór:} GRU/LSTM nie jako model główny, lecz jako \emph{baseline} w~PyTorch do oceny kompromisu dokładność--czas.

  \item[\cite{chung2014gru} --- Chung i~in., 2014]
  GRU upraszcza LSTM przy zbliżonej jakości. \textbf{Wniosek:} jeśli testować rekurencję, GRU jest tańszy. \textbf{Wybór:} dwuwarstwowe GRU jako punkt odniesienia.

  \item[\cite{nie2023patchtst} --- Nie i~in., ICLR 2023, PatchTST]
  Transformer na ``łatki'' szeregu czasowego, kanały niezależnie. Dobrze modeluje długi horyzont, ale jest pamięciożerny. \textbf{Wniosek:} dla $T=1500$ pełna uwaga jest kosztowna; PatchTST jest realistycznym SOTA do porównania, nie do własnej implementacji od zera. \textbf{Wybór:} PatchTST z~HuggingFace wyłącznie jako referencja, nie jako rdzeń biblioteki.
\end{description}

\subsection{Reprezentacje wejścia: raw, magnitude, hybrid}

\begin{description}[style=nextline, leftmargin=0.6em, font=\normalfont]
  \item[\cite{figo2010preprocessing} --- Figo i~in., 2010]
  Przegląd preprocessingu IMU: osie surowe vs.\ cechy niezmiennicze (m.in.\ norma wektora przyspieszenia). \textbf{Wniosek:} magnitude $m_t=\lVert (x_t,y_t,z_t)\rVert_2$ jest tania i~odporna na obrót obudowy, ale gubi kierunek. \textbf{Wybór:} osobny wariant \textit{mag} (norma acc + norma gyro).

  \item[\cite{yurtman2017orientation} --- Yurtman, Barshan, Sensors 2017]
  Pokazują, że zmiana orientacji czujnika silnie psuje klasyfikatory na surowych osiach; transformaty niezmiennicze odzyskują część jakości. \textbf{Wniosek:} sam \textit{raw} może być kruchy, sam \textit{mag} --- zbyt ubogi dla klas różniących się orientacją ciała (stanie vs.\ leżenie). \textbf{Wybór:} wariant \textit{hybrid}: 6 osi + 2 moduły, aby sieć sama ważyła kierunek i~energię ruchu.

  \item[\cite{shoaiB2014fusion} --- Shoaib i~in., 2014]
  Fuzja akcelerometru i~żyroskopu poprawia HAR względem jednego czujnika. \textbf{Wniosek:} nie redukujemy wejścia do samego acc. \textbf{Wybór:} zawsze oba czujniki, we wszystkich trzech trybach.
\end{description}

\subsection{Elementy uczenia głębokiego (do samodzielnej implementacji)}

\begin{description}[style=nextline, leftmargin=0.6em, font=\normalfont]
  \item[\cite{ioffe2015bn} --- Ioffe, Szegedy, Batch Normalization, 2015]
  Stabilizuje rozkład aktywacji, pozwala na wyższe $\eta$. \textbf{Wniosek:} BN jest standardem w~CNN i~musi mieć tryb train/eval (statystyki kroczące). \textbf{Wybór:} BN po każdej konwolucji i~po pierwszej warstwie gęstej.

  \item[\cite{kingma2015adam} --- Kingma, Ba, Adam, 2015]
  Adaptacyjne momenty; domyślne $\beta_1=0{,}9$, $\beta_2=0{,}999$. \textbf{Wniosek:} Adam jest mniej wrażliwy na skalę gradientu niż SGD -- istotne przy własnej, początkowo niestabilnej implementacji. \textbf{Wybór:} Adam jako optymalizator domyślny, SGD tylko do testów poprawności.

  \item[\cite{he2015delving} --- He i~in., inicjalizacja, 2015]
  Skalowanie $\sqrt{2/n_{\mathrm{in}}}$ dla ReLU/wariantów. \textbf{Wniosek:} zła inicjalizacja psuje Conv1D przy dużych jądrach. \textbf{Wybór:} He dla Dense i~Conv1D.

  \item[\cite{srivastava2014dropout} --- Srivastava i~in., 2014; inverted dropout]
  \textbf{Wybór:} dropout z~odwrotnym skalowaniem w~treningu, aby inferencja nie wymagała mnożenia wag.

  \item[\cite{clevert2016elu,maas2013lrelu} --- ELU / Leaky ReLU]
  Łagodzą problem martwych neuronów ReLU. \textbf{Wybór:} LeakyReLU we wczesnych blokach, ELU głębiej i~w~głowie klasyfikatora -- kompromis między tanim nieliniowym mapowaniem a~płynnym gradientem.

  \item[\cite{lin2014gap} --- Lin i~in., Network-in-Network / global average pooling]
  GAP zastępuje spłaszczenie, zmniejsza overfitting i~uniezależnia klasyfikator od $T$. \textbf{Wniosek:} przy zmiennej długości KU-HAR flatten byłby błędny (padding). \textbf{Wybór:} \emph{concat pool} (średnia + max) z~dzieleniem średniej przez prawdziwą długość sekwencji, nie przez $1500$.
\end{description}

\subsection{Obliczenia na GPU i~istniejące biblioteki}

\begin{description}[style=nextline, leftmargin=0.6em, font=\normalfont]
  \item[\cite{okuta2017cupy} --- Okuta i~in., CuPy, LearningSys @ NIPS 2017]
  CuPy odwzorowuje API NumPy na CUDA (cuBLAS i~pokrewne). Wymiana \texttt{import numpy as np} na \texttt{cupy} przenosi obliczenia na GPU. \textbf{Wniosek:} własna biblioteka może być przenośna CPU/GPU bez kerneli CUDA pisanych ręcznie. \textbf{Wybór:} warstwa \texttt{backend.xp} jako alias tablic; to główny wkład HPC projektu.

  \item[\cite{harris2007gpu} --- Harris, GPU GEMM / redukcja; model SIMT]
  Efektywność GPU wymaga regularnego, wektoryzowanego dostępu do pamięci i~dużych paczek. \textbf{Wniosek:} pętle Python nad czasem $T$ zniszczą speedup. \textbf{Wybór:} Conv1D przez \texttt{as\_strided} + \texttt{tensordot} (im2col), batch $\geq 128$, \texttt{float32}.

  \item[\cite{chetlur2014cudnn} --- Chetlur i~in., cuDNN]
  Dojrzałe frameworki korzystają ze stopionych kerneli splotu. \textbf{Wniosek:} własna Conv1D będzie wolniejsza od PyTorch -- to oczekiwane i~warto zmierzyć, nie ukrywać. \textbf{Wybór:} identyczna architektura w~PyTorch jako górne ograniczenie wydajności.

  \item[\cite{paszke2019pytorch} --- Paszke i~in., PyTorch]
  Dynamiczny graf, autograd, przezroczyste GPU. \textbf{Wniosek:} PyTorch jest właściwym odniesieniem poprawności gradientów i~czasu, ale przesłania mechanikę wstecznej propagacji -- stąd własna implementacja na potrzeby dydaktyczne przedmiotu.

  \item[\cite{van2011numpy} --- van der Walt i~in., NumPy]
  \textbf{Wybór:} NumPy jako backend referencyjny (łatwy debug, zgodność z~testami jednostkowymi na CPU).
\end{description}

\section{Uzasadnienie wyboru metod}

Na podstawie powyższego przeglądu przyjmujemy następujące decyzje \textbf{projektowe} (do weryfikacji empirycznej w~dalszej części przedmiotu).

\begin{enumerate}[leftmargin=1.4em, itemsep=0.4em]
  \item \textbf{Własna biblioteka, nie sam PyTorch.}
  Przedmiot wymaga świadomego sterowania obliczeniami (układ pamięci, wybór urządzenia, koszt warstwy), a~nie wyłącznie wywołania \texttt{nn.Module}. PyTorch pozostaje narzędziem kontrolnym.

  \item \textbf{Conv1D zamiast GRU jako model główny.}
  Splot jest równoległy w~czasie i~w~batchu, więc lepiej saturuje GPU~\cite{yang2015cnn,harris2007gpu}. GRU testujemy później, spodziewając się gorszego wykorzystania akceleratora~\cite{hammerla2016deep}.

  \item \textbf{Trzy wejścia: raw / magnitude / hybrid.}
  Literatura pokazuje konflikt: osie dają kierunek, moduł daje niezmienniczość na obrót~\cite{figo2010preprocessing,yurtman2017orientation}. Porównanie trzech trybów \emph{przy identycznej sieci} pozwoli oddzielić wpływ cech od wpływu architektury.

  \item \textbf{Maskowany concat-pooling zamiast Flatten.}
  KU-HAR ma pole długości i~zero-padding~\cite{sikder2021kuhar}; średnia po całym $T$ fałszowałaby sygnał. Max wnosi piki (skok, uderzenie), średnia -- tło aktywności.

  \item \textbf{Adam + BN + Dropout + early stopping.}
  Zestaw minimalny, który w~literaturze CNN-HAR stabilizuje trening bez ręcznego strojenia SGD~\cite{ioffe2015bn,kingma2015adam,ronao2016cnn}.

  \item \textbf{CuPy jako warstwa superkomputerowa.}
  Zamiast pisać CUDA C++, odwzorowujemy NumPy na GPU~\cite{okuta2017cupy} i~planujemy zmierzyć speedup względem CPU na tej samej ścieżce kodu. To jest mierzalny cel HPC: nie ``najlepszy model na świecie'', lecz \emph{przenośna, wektoryzowana} implementacja, której wąskie gardło da się porównać z~cuDNN.
\end{enumerate}

Planowana architektura docelowa (szkic):
\begin{center}
\small
\begin{tabular}{@{}llll@{}}
\toprule
Blok & Jednostki & Kernel / stride / pad & Dalej \\
\midrule
Conv1D & 32  & $31/1/15$ & BN, LeakyReLU \\
Conv1D & 64  & $15/2/7$  & BN, LeakyReLU, Dropout $0{,}2$ \\
Conv1D & 128 & $7/2/3$   & BN, ELU \\
GlobalConcatPool1D & $256$ cech & maska długości & --- \\
Dense & 128 & --- & BN, ELU, Dropout $0{,}5$ \\
Dense + Softmax & 18 klas & --- & entropia krzyżowa \\
\bottomrule
\end{tabular}
\end{center}

\section{Podsumowanie przeglądu}

Literatura HAR wskazuje CNN 1D na surowym IMU jako dojrzały, obliczeniowo korzystny punkt startu, a~literatura GPU -- CuPy jako najkrótszą drogę do akceleracji kodu tablicowego bez rezygnacji z~NumPy. Luka, którą projekt zamierza wypełnić, nie polega na nowym algorytmie SOTA, lecz na \textbf{spójnym połączeniu}: własna, przezroczysta implementacja wstecznej propagacji, trzy świadome reprezentacje IMU oraz pomiar, ile wydajności traci się względem PyTorch/cuDNN, zyskując pełną kontrolę nad jądrem obliczeniowym. Na tej podstawie przechodzimy do implementacji.

\begin{thebibliography}{99}

\bibitem{sikder2021kuhar}
N.~Sikder, A.-A.~Nahid,
\textit{KU-HAR: An open dataset for heterogeneous human activity recognition},
Pattern Recognition Letters, 146, 46--54, 2021.

\bibitem{anguita2013har}
D.~Anguita, A.~Ghio, L.~Oneto, X.~Parra, J.~L.~Reyes-Ortiz,
\textit{A public domain dataset for human activity recognition using smartphones},
ESANN, 2013.

\bibitem{wang2019survey}
J.~Wang, Y.~Chen, S.~Hao, X.~Peng, L.~Hu,
\textit{Deep learning for sensor-based activity recognition: A survey},
Pattern Recognition Letters, 119, 3--11, 2019.

\bibitem{yang2015cnn}
J.~B.~Yang, M.~N.~Nguyen, P.~P.~San, X.~L.~Li, S.~Krishnaswamy,
\textit{Deep convolutional neural networks on multichannel time series for human activity recognition},
IJCAI, 2015.

\bibitem{ronao2016cnn}
C.~A.~Ronao, S.-B.~Cho,
\textit{Human activity recognition with smartphone sensors using deep learning architectures},
Expert Systems with Applications, 59, 235--246, 2016.

\bibitem{hammerla2016deep}
N.~Y.~Hammerla, S.~Halloran, T.~Plötz,
\textit{Deep, convolutional, and recurrent models for human activity recognition using wearables},
IJCAI, 2016.

\bibitem{chung2014gru}
J.~Chung, C.~Gulcehre, K.~Cho, Y.~Bengio,
\textit{Empirical evaluation of gated recurrent neural networks on sequence modeling},
arXiv:1412.3555, 2014.

\bibitem{nie2023patchtst}
Y.~Nie, N.~H.~Nguyen, P.~Sinthong, J.~Kalagnanam,
\textit{A time series is worth 64 words: Long-term forecasting with transformers},
ICLR, 2023.

\bibitem{figo2010preprocessing}
D.~Figo, P.~C.~Diniz, D.~R.~Ferreira, J.~M.~P.~Cardoso,
\textit{Preprocessing techniques for context recognition from accelerometer data},
Personal and Ubiquitous Computing, 14(7), 645--662, 2010.

\bibitem{yurtman2017orientation}
A.~Yurtman, B.~Barshan,
\textit{Activity recognition invariant to sensor orientation with wearable motion sensors},
Sensors, 17(8), 1838, 2017.

\bibitem{shoaiB2014fusion}
M.~Shoaib, S.~Bosch, O.~D.~Incel, H.~Scholten, P.~J.~M.~Havinga,
\textit{Fusion of smartphone motion sensors for physical activity recognition},
Sensors, 14(6), 10146--10176, 2014.

\bibitem{ioffe2015bn}
S.~Ioffe, C.~Szegedy,
\textit{Batch normalization: Accelerating deep network training by reducing internal covariate shift},
ICML, 2015.

\bibitem{kingma2015adam}
D.~P.~Kingma, J.~Ba,
\textit{Adam: A method for stochastic optimization},
ICLR, 2015.

\bibitem{he2015delving}
K.~He, X.~Zhang, S.~Ren, J.~Sun,
\textit{Delving deep into rectifiers: Surpassing human-level performance on ImageNet classification},
ICCV, 2015.

\bibitem{srivastava2014dropout}
N.~Srivastava, G.~Hinton, A.~Krizhevsky, I.~Sutskever, R.~Salakhutdinov,
\textit{Dropout: A simple way to prevent neural networks from overfitting},
JMLR, 15, 1929--1958, 2014.

\bibitem{clevert2016elu}
D.-A.~Clevert, T.~Unterthiner, S.~Hochreiter,
\textit{Fast and accurate deep network learning by exponential linear units (ELUs)},
ICLR, 2016.

\bibitem{maas2013lrelu}
A.~L.~Maas, A.~Y.~Hannun, A.~Y.~Ng,
\textit{Rectifier nonlinearities improve neural network acoustic models},
ICML Workshop, 2013.

\bibitem{lin2014gap}
M.~Lin, Q.~Chen, S.~Yan,
\textit{Network in network},
ICLR, 2014.

\bibitem{okuta2017cupy}
R.~Okuta, Y.~Unno, D.~Nishino, S.~Hido, C.~Loomis,
\textit{CuPy: A NumPy-compatible library for NVIDIA GPU calculations},
LearningSys Workshop, NIPS, 2017.

\bibitem{harris2007gpu}
M.~Harris,
\textit{Optimizing parallel reduction in CUDA},
NVIDIA Developer Technology, 2007.

\bibitem{chetlur2014cudnn}
S.~Chetlur i~in.,
\textit{cuDNN: Efficient primitives for deep learning},
arXiv:1410.0759, 2014.

\bibitem{paszke2019pytorch}
A.~Paszke i~in.,
\textit{PyTorch: An imperative style, high-performance deep learning library},
NeurIPS, 2019.

\bibitem{van2011numpy}
S.~van der Walt, S.~C.~Colbert, G.~Varoquaux,
\textit{The NumPy array: A structure for efficient numerical computation},
Computing in Science \& Engineering, 13(2), 22--30, 2011.

\end{thebibliography}

\end{document}
